\section{Preface}
  \khnote{test}
  \hhnote{1st draft done}

  \noindent
  So you want to prove post-quantum security in the random oracle model.

  Proofs of security in the ROM are often conceptually
    simple, and yet regularly achieve strong security claims with
    near-tightness for highly efficient constructions---which is as close as
    one gets to a holy trinity for provable security. In contrast, proofs in the
    Quantum Random Oracle Model are often hard to parse without specialized
    knowledge in quantum information theory, and infamously lossy.

  And yet, if we want our proofs of security to have merit in a post-quantum
    world, it seems that our choice is between the Standard Model and the
    QROM. This is because random oracles model hash functions, which are
    publicly available algorithms that anyone can run locally on their machine.
    Therefore, there is nothing stopping someone with a quantum computer to
    boot it up and run it on a superposition of inputs. If we want to
    accurately model quantum attackers, then, our model must allow for such
    strategies. 

  In the standard model, this is not a problem at all,
    since then the hash functions in question is provided to the attacker
    as-is. But if QROM proofs are lossy, then Standard Model constructions
    that achieve the same goals as their ROM counterparts are often highly
    inefficient in comparison (and, as every cryptographer knows, efficiency
    tends to trump security in the real world).

  Therefore, as we inch ever closer to the looming threat of
    cryptographically-relevant quantum computers, it is important that the
    barrier for proving security in the QROM is lowered as much as
    possible, so that the community may together bring our transition to a
    post-quantum world into effect, in a safe, efficient, and (provably) secure
    manner.

  \paragraph{ROM Techniques Don't Transfer \ldots}
    One of the most useful facts about security reductions in the random oracle model
      is that anyone can efficiently \emph{simulate} a random oracle, simply by
      drawing the random outputs one-at-a-time and storing them in an
      appropriate list. This technique, known as \emph{lazy sampling}, allows
      the reduction to not only implement the oracle efficiently, but also to
      see what input values the adversary is interested in, \emph{and}
      adaptively reprogram the oracle on inputs of its choice.
      Indeed, much of the strength of the random oracle model lies in exactly
      these capabilities.

    None of this transfers directly to the case of quantum-accessible oracles.
      To see why, consider an adversary that queries the oracle on a
      superposition of all possible inputs at once: the simulator would have to
      store an exponentially large list of input values in order process just
      this single query!

  \paragraph{\ldots But the Concepts Do.}
    And yet, the past decade has seen the development of techniques that
      recover most of these capabilities, albeit at a loss. Prime among these
      techniques are the various One-Way to Hiding (O2H) lemmas, which allow
      for careful reprogramming of the quantum random oracle (as well as
      limited extraction of the input values); Quantum Online-Extractability
      (QOE), which as the name suggests allows for on-the-fly extraction of oracle inputs;
      and the Compressed Random Oracle (CRO) technique, which allows one to
      efficiently simulate a quantum random oracle, in direct
      analogy to classical ``lazy sampling''.

  \subsubsection{Outline.}
    %We start %in \autoref{part:prelims} 
      We begin in Part 1 by laying the foundations, with
      definitions for basic crypto concepts in \autoref{sec:basics}, and a
      primer on quantum computing concepts in \autoref{sec:primer}.
      In \autoref{sec:pseudocode}, we introduce a novel pseudocode notation,
      which will be helpful in reasoning about oracles and reductions.
      Then in \autoref{sec:rom}, we finally introduce our main objects of study,
      the quantum random oracles. Through illustrative examples we will be
      building intuition, aiming to get a ``feel'' for how they work.

    \hhnote{update the following once the v0.1 structure is settled}

    \autoref{part:classical} is all about the classical case. Here, we show
      several examples of security proofs in the random oracle model, all of
      which are conceptually simple and yield (relatively) tight security.
      The examples are lifted from the security proofs of two important
      transformations, the Fujisaki-Okamoto (FO) transformation (which transforms
      a weakly-secure public-key encryption scheme into a strongly-secure key
      encapsulation mechanism), and the Fiat-Shamir (FS) transformation (which
      transforms an interactive zero-knowledge protocol into a non-interactive
      signature scheme). 

    We meet these examples again in \autoref{part:tools}, where we use
      them to showcase the tools as we introduce them by showing that security
      holds also in the QROM, though usually with a security loss compared to the
      classical case, with the particular security loss depending on the tool used.
      \autoref{sec:o2h} is all about the One-way To Hiding (O2H) lemma, and its many variants: 
      how to use it to reprogram quantum random oracles, when to employ which
      version, and the impact of each version on the tightness of the proof. 
      Our running example here
      will be the second half of the FO transformation, known as the U-transform,
      and showing that if we start with a one-way secure (OW-CPA) deterministic PKE, we
      get an indistinguishable KEM secure against passive attacks (IND-CPA). 
      In \autoref{sec:extract}, we introduce the technique of
      Quantum Online-Extraction (QOE), and show how it can be used to prove
      that the U-transform \emph{also} upgrades to security against active attacks
      (IND-CCA).
      In \autoref{sec:adaptive}, we show how reprogramming can be made
      \emph{adaptive}, meaning that the reprogramming may now happen
      on-the-fly, on points that may (in full or in part) depend on
      adversarially-chosen input values. We use this to show
      how to go from a passively-secure Non-Interactive Zero-Knowledge (NIZK)
      protocol (such as constructed from a sigma protocol via the FS
      transformation) to an actively-secure signature scheme (EUF-CMA).
      Then, in \autoref{sec:compressed}, we delve into the details of 
      Compressed Random Oracles (CROs), and show how one might implement the
      technique to provide an efficient yet exact simulation of a quantum
      random oracle.
      This incredibly powerful technique takes some effort to wrap one's head
      around, and so we dedicate significant space to give gentle and
      intuitive explanations for all its moving parts, with several examples
      showing how queries are processed.
    \ifdraftversion
      Finally, in \autoref{sec:rewind}, we show how to lift rewinding arguments
      from the ROM to the QROM, with the central example being how the FS
      transformation transforms a kind of interactive zero-knowledge protocol known
      as a sigma protocol, to a NIZK.
      \hhnote{leave for version 2}
    \fi

    \ifdraftversion
    \autoref{part:proofs} is all about \emph{constructing the tools
      themselves}, which requires a deeper dive into the world of quantum
      information theory.
      \hhnote{and so on and so forth; leave for version 3}
    \fi

    The appendices contain extra material, including some more formal
      definitions in \autoref{sec:pseudocode}, and solutions to select
      exercises in \autoref{sec:solutions}.

    \hhnote{fix ``Sect.'' vs ``App.''}
